\documentclass[12pt, a4paper]{article}

% --- Packages ---
\usepackage[utf8]{inputenc}
\usepackage[T1]{fontenc}
\usepackage[french]{babel}
\usepackage{graphicx}
\usepackage{amsmath, amssymb}
\usepackage{booktabs}
\usepackage{hyperref}
\usepackage{geometry}
\usepackage{float}
\usepackage{caption}
\usepackage{subcaption}
\usepackage{xcolor}
\usepackage{listings}
\usepackage{fancyhdr}
\usepackage{enumitem}
\usepackage{array}
\usepackage{multirow}
\usepackage{colortbl}
\usepackage{tabularx}
\usepackage{tcolorbox}
\tcbuselibrary{skins}

\geometry{margin=2.5cm}

% --- En-tete / Pied de page ---
\pagestyle{fancy}
\fancyhf{}
\fancyhead[L]{\small Hackathon -- Prediction de series temporelles financieres}
\fancyhead[R]{\small\thepage}
\fancyfoot[C]{\small Louis Cluzel -- M1 IA -- Universite Clermont Auvergne -- 2026}

% --- Couleurs ---
\definecolor{vertclair}{RGB}{198,239,206}
\definecolor{rougeclair}{RGB}{255,199,206}
\definecolor{orangeclair}{RGB}{255,235,156}
\definecolor{bleutitre}{RGB}{21,101,192}
\definecolor{gristexte}{RGB}{80,80,80}

% --- Style listings Python ---
\lstset{
  language=Python,
  basicstyle=\ttfamily\footnotesize,
  keywordstyle=\color{bleutitre}\bfseries,
  commentstyle=\color{gray},
  stringstyle=\color{red!70!black},
  numbers=left,
  numberstyle=\tiny\color{gray},
  frame=single,
  breaklines=true,
  captionpos=b,
  backgroundcolor=\color{gray!5}
}

% --- Boite mise en avant ---
\newtcolorbox{resultbox}{
  enhanced, colback=green!6, colframe=green!50!black,
  boxrule=1pt, arc=4pt, left=6pt, right=6pt, top=4pt, bottom=4pt
}
\newtcolorbox{formulebox}{
  enhanced, colback=blue!5, colframe=bleutitre,
  boxrule=1pt, arc=4pt, left=6pt, right=6pt, top=4pt, bottom=4pt
}

% --- Page de titre ---
\title{
  \vspace{-0.5cm}
  \includegraphics[width=4cm]{figures/AAPL_indicators_overview.png}\\[0.4cm]
  \textbf{\Large Hackathon : Prediction de series temporelles}\\[0.3cm]
  \textbf{\Large appliquee aux indicateurs financiers}\\[0.5cm]
  \large Pipeline multifeature inter-actifs (AAPL, MSFT, NVDA)\\[0.2cm]
  \normalsize Comparaison LSTM~$\cdot$~GRU~$\cdot$~Transformer~$\cdot$~Ensemble~$\cdot$~Baseline\\
  sur MACD(12,26,9) et Stochastique(14,3,5)
}
\author{
  \textbf{Louis Cluzel}\\
  \small M1 Intelligence Artificielle\\
  \small Universite Clermont Auvergne
}
\date{Avril 2026}

\begin{document}

\maketitle
\thispagestyle{empty}
\vspace{0.5cm}

\begin{abstract}
\noindent
Ce rapport presente un pipeline unifie de prediction de series temporelles applique aux indicateurs financiers MACD(12,26,9) et Stochastique \%K(14,3,5). Il evalue la robustesse de cinq architectures de prediction (Random Walk, ARIMA, LSTM, GRU et Transformer) sur un horizon de 15 pas de temps, testees de maniere homogene sur trois actifs de nature differente~: Apple (intraday 5~min, 94~334~barres), Microsoft et Nvidia (journalier, $\sim$1~250~jours).

\noindent
L'evaluation repose sur le \textbf{score hackathon} $S = (T_\text{MACD} + \max(0,\rho_\text{MACD}) + T_\text{Stoch} + \max(0,\rho_\text{Stoch}))/4$, avec un objectif $S \geq 0{,}55$. Les resultats revelent une forte dependance du modele optimal au regime de marche de l'actif~: le GRU domine sur AAPL ($S=0{,}79$), le LSTM l'emporte sur NVDA ($S=0{,}59$), et le Transformer se distingue sur MSFT par sa precision stochastique.
\end{abstract}

\tableofcontents
\newpage

% =========================================================================
\section{Introduction}
% =========================================================================

La prediction des indicateurs techniques financiers constitue un probleme canonique de prediction de series temporelles non stationnaires. Ces indicateurs, derives de prix et de volumes, presentent des dynamiques non lineaires et une sensibilite aux regimes de marche qui rendent leur anticipation particulierement ardue.

L'objectif de ce projet est de construire un pipeline modulaire permettant de predire l'evolution sur \textbf{15 pas de temps futurs} de deux indicateurs complementaires~:
\begin{itemize}[noitemsep]
  \item le \textbf{MACD} (12,~26,~9)~: signal tendanciel lisse par difference de moyennes mobiles exponentielles, a dynamique lente et reguliere ;
  \item l'\textbf{oscillateur stochastique \%K} (14,~3,~5)~: indicateur de momentum mesurant la position du cours dans son range recent, inherentement cyclique et bruite.
\end{itemize}

L'originalite de l'approche reside dans l'evaluation systematique et homogene du pipeline sur trois actifs de profils financiers distincts (AAPL, MSFT, NVDA), permettant d'identifier les forces et faiblesses de chaque architecture face a differents regimes de marche.

% =========================================================================
\section{Donnees et Pre-traitement}
% =========================================================================

\subsection{Source des donnees}

L'approche multi-actifs couvre volontairement des periodes et granularites differentes afin de stresser le pipeline~:

\begin{table}[H]
\centering
\caption{Description des jeux de donnees utilises}
\begin{tabular}{llccc}
\toprule
\textbf{Actif} & \textbf{Secteur} & \textbf{Granularite} & \textbf{Barres} & \textbf{Periode} \\
\midrule
AAPL (Apple)      & Tech mature   & Intraday 5~min & 94~334 & 2010--2024 \\
MSFT (Microsoft)  & Tech mature   & Journalier     & 1~236  & 2021--2026 \\
NVDA (Nvidia)     & Semi-conducteurs & Journalier  & 1~256  & 2021--2026 \\
\bottomrule
\end{tabular}
\end{table}

Ce choix permet de confronter les modeles a~: (i)~un historique profond et regulier (AAPL), (ii)~un actif stable de grande capitalisation (MSFT), et (iii)~un titre extremement volatil en phase de bulle technologique (NVDA).

\subsection{Feature engineering}

Chaque pas de temps $t$ est represente par un vecteur de 5 composantes normalisees~:
\begin{table}[H]
\centering
\caption{Features calculees pour chaque pas de temps}
\begin{tabular}{lll}
\toprule
\textbf{Feature} & \textbf{Description} & \textbf{Role} \\
\midrule
\texttt{value}  & Valeur normalisee MinMax [0,1] & Signal courant \\
\texttt{return} & Variation relative & Dynamique instantanee \\
\texttt{ma10}   & Moyenne mobile 10 periodes & Tendance court terme \\
\texttt{ma20}   & Moyenne mobile 20 periodes & Tendance moyen terme \\
\texttt{std10}  & Ecart-type glissant 10 periodes & Volatilite locale \\
\bottomrule
\end{tabular}
\end{table}

Ces features sont injectees dans des fenetres glissantes de lookback $L = 120$ pas, produisant des sequences supervisees $X_i \in \mathbb{R}^{120 \times 5}$ associees a un horizon de prediction $y_i \in \mathbb{R}^{15}$. La normalisation est ajustee uniquement sur les donnees d'entrainement pour prevenir toute fuite d'information.

% =========================================================================
\section{Methodologie}
% =========================================================================

Le pipeline evalue six approches de complexite croissante~:

\begin{itemize}
  \item \textbf{Random Walk with Drift (Baseline)}~: extrapolation lineaire deterministe sur la derive historique. Sert de borne inferieure.
  \item \textbf{ARIMA(5,1,2)}~: modele autoregressif classique en \textit{rolling forecast}, reajuste a chaque pas.
  \item \textbf{LSTM}~: reseau recurrent a memoire longue gerant un etat de cellule $C_t$ separe de l'etat cache $h_t$ via des portes d'oubli, d'entree et de sortie. Prediction directe multi-pas via une couche $\text{Linear}(64, 15)$.
  \item \textbf{GRU}~: variante simplifiee du LSTM fusionnant $C_t$ et $h_t$ en un seul vecteur avec deux portes. Moins de parametres, convergence plus rapide, meilleure resistance au sur-apprentissage sur de petits echantillons.
  \item \textbf{Transformer encoder}~: mecanisme d'attention globale avec encodage positionnel sinusoidal ($d_\text{model} = 32$, 1 couche, 4 tetes). Prediction sur le dernier token.
  \item \textbf{Ensemble}~: moyenne ponderee des trois modeles DL, poids optimises par Nelder-Mead.
\end{itemize}

\subsection{Score hackathon}

\begin{formulebox}
\begin{equation}
  S = \frac{T_\text{MACD} + \max(0,\, \rho_\text{MACD}) + T_\text{Stoch} + \max(0,\, \rho_\text{Stoch})}{4}
\end{equation}
\textbf{Objectif~: $S \geq 0{,}55$.} Le score de tendance $T$ mesure l'alignement directionnel des variations~; la correlation de Pearson $\rho$ mesure la fidelite de forme. Les correlations negatives sont penalisees (remplacees par 0).
\end{formulebox}

% =========================================================================
\section{Resultats experimentaux}
% =========================================================================

Cette section presente les resultats detailles de chaque actif, suivi d'une analyse croisee approfondie. Les 15 derniers pas de chaque serie constituent le jeu de test~; les modeles DL sont entraines sur les barres precedentes.

% -------------------------------------------------------------------
\subsection{Apple (AAPL) -- Intraday haute frequence}
% -------------------------------------------------------------------

La configuration de reference utilise 6~000 barres d'entrainement et 20 epochs (meilleur compromis identifie sur 3 runs independants).

\begin{table}[H]
\centering
\caption{Resultats detailles -- AAPL (Intraday 5~min, 6~000 barres, 20 epochs)}
\begin{tabular}{clccccc}
\toprule
\textbf{Rang} & \textbf{Modele} & $T_\text{MACD}$ & $\rho_\text{MACD}$ & $T_\text{Stoch}$ & $\rho_\text{Stoch}$ & \textbf{Score} \\
\midrule
\rowcolor{vertclair}
1 & \textbf{GRU}   & 78,6\% & 0,9657 & 64,3\% & 0,7690 & \textbf{0,7908} \\
2 & LSTM           & 85,7\% & 0,9516 & 57,1\% & 0,4345 & 0,7037 \\
3 & Ensemble       & 78,6\% & 0,9461 & 50,0\% & 0,1058 & 0,5844 \\
4 & Transformer    & 50,0\% & 0,6850 & 28,6\% & $-$0,2671 & 0,3677 \\
\rowcolor{rougeclair}
5 & Baseline       & 21,4\% & $-$0,9658 & 57,1\% & 0,3010 & 0,2717 \\
\bottomrule
\end{tabular}
\end{table}

\begin{figure}[H]
  \centering
  \begin{subfigure}[b]{0.48\textwidth}
    \includegraphics[width=\textwidth]{figures/fig_pred_macd.png}
    \caption{MACD Line -- 15 pas}
    \label{fig:aapl_macd}
  \end{subfigure}
  \hfill
  \begin{subfigure}[b]{0.48\textwidth}
    \includegraphics[width=\textwidth]{figures/fig_pred_stoch.png}
    \caption{Stochastique \%K -- 15 pas}
    \label{fig:aapl_stoch}
  \end{subfigure}
  \caption{Predictions LSTM et GRU vs valeurs reelles (AAPL)}
  \label{fig:predictions_aapl}
\end{figure}

\textbf{Analyse detaillee.}
Le GRU domine nettement avec un score de \textbf{0,7908}, le plus eleve de l'ensemble des experimentations. Plusieurs observations meritent attention~:

\begin{itemize}
  \item \textbf{MACD~: prediction fiable pour tous les modeles DL.} Les correlations MACD des trois reseaux profonds depassent 0,68 (GRU~: 0,97, LSTM~: 0,95, Transformer~: 0,69). Le MACD, construit sur des moyennes mobiles exponentielles lentes (periodes 12 et 26), genere une serie intrinsequement lisse dont la tendance a court terme est previsible. L'ecart entre recurrents ($\rho > 0{,}94$) et Transformer ($\rho = 0{,}69$) suggere que la memoire sequentielle ordonnee des cellules recurrentes capte mieux l'inertie du MACD que l'attention globale.

  \item \textbf{Stochastique~: facteur discriminant du classement.} La differentiation decisive entre GRU ($\rho_\text{Stoch} = 0{,}77$) et LSTM ($\rho_\text{Stoch} = 0{,}43$) se joue sur le stochastique. Le GRU, avec moins de parametres, explore un espace de solutions plus reduit et converge plus efficacement vers un optimum adapte aux oscillations rapides ($\sim$14~barres de cycle) du \%K. Le LSTM, sur-parametre pour cette tache, semble stocker de la memoire a long terme non pertinente qui bruite les gradients.

  \item \textbf{Transformer en difficulte.} Le score de 0,37 (en dessous de la baseline stochastique) s'explique par la triple contrainte~: $d_\text{model} = 32$ impose par le CPU, attention globale inadaptee aux patterns locaux de 120~barres, et volume de sequences ($\sim$5~800) insuffisant pour ce type d'architecture.
\end{itemize}

% -------------------------------------------------------------------
\subsection{Microsoft (MSFT) -- Journalier, actif stable}
% -------------------------------------------------------------------

\begin{table}[H]
\centering
\caption{Resultats detailles -- MSFT (Journalier, $\sim$1~200 barres)}
\begin{tabular}{clccccccc}
\toprule
\textbf{Rang} & \textbf{Modele} & $T_\text{MACD}$ & $\rho_\text{MACD}$ & \textbf{MAE\textsubscript{MACD}} & $T_\text{Stoch}$ & $\rho_\text{Stoch}$ & \textbf{MAE\textsubscript{Stoch}} & \textbf{Score} \\
\midrule
\rowcolor{orangeclair}
1 & \textbf{Transformer} & 50,0\% & $-$0,33 & 19,9 & 57,1\% & 0,81 & 16,5 & \textbf{0,4695} \\
2 & GRU       & 42,9\% & $-$0,05 & 13,5 & 57,1\% & 0,85 & 18,1 & 0,4614 \\
3 & Ensemble  & 42,9\% & $-$0,42 & 13,5 & 57,1\% & 0,83 & 17,8 & 0,4580 \\
4 & LSTM      & 35,7\% & $-$0,78 & 33,4 & 57,1\% & 0,83 & 19,8 & 0,4387 \\
5 & ARIMA     &  7,1\% & $-$0,88 &  8,7 & 35,7\% & 0,11 & 52,0 & 0,1356 \\
\rowcolor{rougeclair}
6 & Baseline  &  7,1\% & $-$0,97 &  7,7 & 42,9\% & $-$0,70 & 46,7 & 0,1250 \\
\bottomrule
\end{tabular}
\end{table}

\begin{figure}[H]
  \centering
  \begin{subfigure}[b]{0.48\textwidth}
    \includegraphics[width=\textwidth]{figures/MSFT_MACD_Line_predictions.png}
    \caption{MACD Line -- 15 pas}
    \label{fig:msft_macd}
  \end{subfigure}
  \hfill
  \begin{subfigure}[b]{0.48\textwidth}
    \includegraphics[width=\textwidth]{figures/MSFT_Stochastique_pctK_predictions.png}
    \caption{Stochastique \%K -- 15 pas}
    \label{fig:msft_stoch}
  \end{subfigure}
  \caption{Predictions des modeles DL et ARIMA (MSFT)}
  \label{fig:predictions_msft}
\end{figure}

\textbf{Analyse detaillee.}
Le scenario MSFT revele un comportement radicalement different de celui d'AAPL~:

\begin{itemize}
  \item \textbf{Effondrement generalise du MACD.} C'est le fait marquant~: \emph{tous} les modeles, y compris les DL, affichent des correlations MACD negatives (de $-0{,}05$ pour le GRU a $-0{,}97$ pour la Baseline). Le MACD de MSFT se trouvait dans une phase de retournement de tendance au moment du test~: les modeles, entraines sur un historique ou le MACD evoluait dans une direction, ont extrapole cette tendance alors que le signal reel s'est inverse. Ce phenomene illustre la fragilite des modeles sur des points d'inflexion. Le fait que le GRU ait la correlation la moins negative ($-0{,}05$) suggere qu'il a mieux detecte le debut du retournement.

  \item \textbf{Stochastique~: unanimite remarquable.} En contraste total avec le MACD, les quatre modeles DL convergent vers des correlations stochastiques elevees ($\rho \in [0{,}81\text{--}0{,}85]$)~: le stochastique journalier de MSFT presentait un signal oscillant regulier bien capture par tous les reseaux. Le GRU obtient meme le meilleur $\rho_\text{Stoch}$ (0,85) mais perd le classement global a cause de sa $T_\text{MACD}$ plus basse. Le Transformer, malgre son echec sur le MACD, a la meilleure tendance directionnelle MACD (50\%) des DL, ce qui suffit pour prendre la premiere place.

  \item \textbf{Ecart MAE revelateur.} Le LSTM affiche une MAE MACD de 33,4 contre 13,5 pour le GRU~: il a sur-extrapole la tendance descendante du MACD, projetant des valeurs beaucoup trop eloignees de la realite. Un exces de memoire long-terme a agi comme un biais de confirmation sur la tendance historique.
\end{itemize}

\begin{resultbox}
\textbf{Observation cle sur MSFT~:} Aucun modele n'atteint le seuil de $S \geq 0{,}55$, mais l'echec est entierement imputable au MACD (retournement non anticipe). Le stochastique est bien predit par tous les modeles DL ($\rho > 0{,}80$).
\end{resultbox}

% -------------------------------------------------------------------
\subsection{Nvidia (NVDA) -- Journalier, haute volatilite}
% -------------------------------------------------------------------

\begin{table}[H]
\centering
\caption{Resultats detailles -- NVDA (Journalier, $\sim$1~250 barres)}
\begin{tabular}{clccccccc}
\toprule
\textbf{Rang} & \textbf{Modele} & $T_\text{MACD}$ & $\rho_\text{MACD}$ & \textbf{MAE\textsubscript{MACD}} & $T_\text{Stoch}$ & $\rho_\text{Stoch}$ & \textbf{MAE\textsubscript{Stoch}} & \textbf{Score} \\
\midrule
\rowcolor{vertclair}
1 & \textbf{LSTM} & 71,4\% & 0,92 & 4,1 & 57,1\% & 0,17 & 21,1 & \textbf{0,5933} \\
2 & GRU           & 78,6\% & 0,90 & 3,8 & 64,3\% & $-$0,37 & 30,2 & 0,5814 \\
3 & Transformer   & 71,4\% & 0,54 & 3,8 & 71,4\% & 0,28 & 28,0 & 0,5620 \\
4 & Ensemble      & 64,3\% & 0,85 & 3,9 & 64,3\% & $-$0,09 & 26,3 & 0,5348 \\
5 & ARIMA         & 50,0\% & 0,65 & 4,6 & 42,9\% & $-$0,83 & 42,9 & 0,3944 \\
\rowcolor{rougeclair}
6 & Baseline      &  7,1\% & $-$0,98 & 5,4 & 35,7\% & $-$0,73 & 41,3 & 0,1071 \\
\bottomrule
\end{tabular}
\end{table}

\textbf{Analyse detaillee.}
NVDA offre le contexte de stress le plus revelateur pour le pipeline~:

\begin{itemize}
  \item \textbf{MACD~: tous les DL reussissent, inversion du classement LSTM/GRU.} Contrairement a AAPL, le LSTM devance le GRU (respectivement $\rho = 0{,}92$ et $0{,}90$). L'explication est structurelle~: le MACD de Nvidia, sur un actif ayant multiplie son cours par plus de 15 en 4 ans, presente des tendances longues et soutenues. La memoire de cellule $C_t$ du LSTM, superflue pour les cycles courts d'AAPL, devient ici un avantage strategique pour capturer l'inertie directionnelle du titre. Les MAE MACD sont tres faibles ($3{,}8$--$4{,}6$), refletant un signal tendanciel coherent sur la macro-periode.

  \item \textbf{Stochastique~: la faiblesse commune.} Paradoxalement, le Transformer obtient la meilleure tendance stochastique ($T = 71{,}4\%$) et une correlation positive ($\rho = 0{,}28$), tandis que le GRU ($\rho = -0{,}37$) et l'Ensemble ($\rho = -0{,}09$) echouent completement. Sur un actif aussi volatil, les oscillations du stochastique entre 0 et 100 sont plus brutales et irregulieres que sur AAPL~: les modeles recurrents, entraines a memoriser des patterns temporels, sont trompes par la rupture de regime entre la phase d'entrainement et la fenetre de test. L'attention globale du Transformer, plus agnostique vis-a-vis de l'ordre sequentiel, tolere mieux cette discontinuite.

  \item \textbf{Competitivite serree.} Les quatre premiers modeles se tiennent dans une marge de 0,06 points ($0{,}53$--$0{,}59$), refletant un compromis constant entre la reussite sur le MACD (recurrents) et la resilience sur le stochastique (Transformer). Le LSTM l'emporte grace a son avance correlative sur le MACD ($+0{,}02$ vs GRU, $+0{,}38$ vs Transformer) qui compense sa faiblesse stochastique.
\end{itemize}

\begin{resultbox}
\textbf{Observation cle sur NVDA~:} Le LSTM et le GRU franchissent tous deux le seuil $S \geq 0{,}55$, confirmant que les reseaux recurrents restent competitifs meme sur des actifs hautement volatils, a condition que le signal tendanciel (MACD) soit suffisamment structure.
\end{resultbox}

% =========================================================================
\section{Analyse transversale}
% =========================================================================

\subsection{Synthese comparative inter-actifs}

\begin{table}[H]
\centering
\caption{Tableau recapitulatif -- Meilleur modele par actif et indicateur}
\begin{tabular}{lccccc}
\toprule
\textbf{Actif} & \textbf{Meilleur global} & \textbf{Score} & \textbf{$\rho$ MACD max} & \textbf{$\rho$ Stoch max} & \textbf{$S \geq 0{,}55$ ?} \\
\midrule
\rowcolor{vertclair}
AAPL & GRU & \textbf{0,7908} & 0,97 (GRU) & 0,77 (GRU) & \checkmark \\
NVDA & LSTM & \textbf{0,5933} & 0,92 (LSTM) & 0,28 (Transf.) & \checkmark \\
\rowcolor{orangeclair}
MSFT & Transformer & \textbf{0,4695} & $-$0,05 (GRU) & 0,85 (GRU) & $\times$ \\
\bottomrule
\end{tabular}
\end{table}

\subsection{Le role determinant de l'indicateur MACD}

L'analyse croisee montre que le MACD est le \textbf{facteur critique de reussite}. Sur AAPL et NVDA, les modeles DL obtiennent des correlations MACD superieures a 0,65, ce qui tire mecaniquement le score global au-dessus de 0,55. Sur MSFT, l'echec est entierement imputable a l'effondrement du MACD (toutes correlations negatives). La raison~: le MACD est un indicateur retarde par construction (il reflete les tendances passees), ce qui le rend previsible \emph{tant que la tendance se maintient}, mais catastrophiquement faux lors de retournements.

\subsection{Asymetrie MACD/Stochastique~: un pattern systematique}

Un resultat remarquable emerge des trois experimentations~: \textbf{les performances sur le MACD et le stochastique sont tendanciellement decorrelees}. Sur AAPL, le GRU domine les deux. Mais sur NVDA, le meilleur modele MACD (LSTM, $\rho = 0{,}92$) n'est que 4\textsuperscript{e} sur le stochastique ($\rho = 0{,}17$), tandis que le Transformer, mediocre sur le MACD ($\rho = 0{,}54$), place la meilleure tendance stochastique ($T = 71{,}4\%$). Sur MSFT, le GRU, dernier en MACD, est premier en stochastique.

Cette asymetrie s'explique par la nature fondamentalement differente des deux indicateurs~: le MACD favorise les modeles a memoire longue (recurrents), tandis que le stochastique, plus local et cyclique, peut etre mieux capte par l'attention positionnelle du Transformer.

\subsection{Impact du volume de donnees}

\begin{table}[H]
\centering
\caption{Correlation entre volume de donnees et performances}
\begin{tabular}{lccc}
\toprule
& \textbf{AAPL} (94k barres) & \textbf{MSFT} (1,2k) & \textbf{NVDA} (1,3k) \\
\midrule
Score max        & 0,7908 & 0,4695 & 0,5933 \\
$\rho_\text{MACD}$ max & 0,97 & $-$0,05 & 0,92 \\
$\rho_\text{Stoch}$ max & 0,77 & 0,85 & 0,28 \\
\bottomrule
\end{tabular}
\end{table}

AAPL, avec 94~334 barres, produit les meilleurs scores absolus. Cependant, NVDA atteint $\rho_\text{MACD} = 0{,}92$ avec seulement 1~256 barres, demontrant que le volume n'est pas le seul facteur~: la \textbf{coherence tendancielle} du signal (forte derive haussiere de NVDA) compense la rarete des donnees. A l'inverse, MSFT, avec un volume comparable a NVDA mais un MACD en retournement, echoue completement sur cet indicateur.

\subsection{Robustesse du modele Ensemble}

Le modele Ensemble se classe systematiquement en 3\textsuperscript{e} ou 4\textsuperscript{e} position (AAPL~: 0,58, MSFT~: 0,46, NVDA~: 0,53). Il n'atteint jamais la premiere place car les poids Nelder-Mead incluent inevitablement une contribution du modele le plus faible (generalement le Transformer sur AAPL/NVDA ou le LSTM sur MSFT), diluant le signal du meilleur modele. L'Ensemble remplit neanmoins son role de filet de securite~: il ne produit jamais le pire score et stabilise les resultats.

% =========================================================================
\section{Discussion}
% =========================================================================

\subsection{Quel modele pour quel regime de marche ?}

Les resultats permettent de degager des recommandations pratiques~:

\begin{table}[H]
\centering
\caption{Recommandation de modele selon le contexte}
\begin{tabular}{lll}
\toprule
\textbf{Contexte} & \textbf{Modele recommande} & \textbf{Justification} \\
\midrule
Trend stable, donnees abondantes & GRU & Convergence rapide, pas de sur-apprentissage \\
Trend fort et soutenu & LSTM & Memoire $C_t$ capture l'inertie longue \\
Peu de donnees, signal bruite & Transformer & Attention globale plus agnostique \\
Incertitude sur le regime & Ensemble & Compromis de robustesse \\
\bottomrule
\end{tabular}
\end{table}

\subsection{Limites identifiees}

\begin{enumerate}
  \item \textbf{Retournements de tendance} : aucun modele n'anticipe correctement un changement de regime (cf. MACD MSFT). Un mecanisme de detection de regime de marche pourrait ameliorer la robustesse.
  \item \textbf{Contrainte CPU} : le Transformer est sous-dimensionne ($d_\text{model} = 32$). Avec $d = 64$ ou 128 et un GPU, les performances seraient probablement meilleures.
  \item \textbf{Granularite heterogene} : la comparaison inter-actifs est limitee par la difference de granularite (intraday vs journalier). Une evaluation sur les memes frequences temporelles renforcerait les conclusions.
  \item \textbf{Absence de seed fixe} : la variabilite inter-runs (observee sur AAPL) introduit un bruit experimental. Un protocole multi-seed ($n \geq 5$) avec intervalles de confiance serait necessaire pour des conclusions statistiquement rigoureuses.
\end{enumerate}

% =========================================================================
\section{Conclusion}
% =========================================================================

Ce projet demontre qu'un pipeline modulaire de prediction d'indicateurs financiers peut s'adapter a des regimes de marche radicalement differents. Les enseignements principaux sont~:

\begin{itemize}
  \item Le \textbf{GRU} est optimal (score de \textbf{0,79}) sur des series regulieres a historique profond (AAPL), grace a sa parcimonie parametrique ;
  \item Le \textbf{LSTM} reprend l'avantage ($S = 0{,}59$) sur les actifs a forte derive directionnelle (NVDA), ou sa memoire longue de cellule capture l'inertie du signal ;
  \item Le \textbf{Transformer}, malgre ses contraintes CPU, se distingue sur la prediction du stochastique dans des contextes a signal bruite (MSFT et NVDA) ;
  \item \textbf{Aucun modele n'est universellement superieur}~: le choix optimal depend intrinsequement de la nature de l'actif et de l'indicateur cible.
\end{itemize}

L'objectif hackathon ($S \geq 0{,}55$) est \textbf{atteint sur AAPL et NVDA}, et le deficit sur MSFT est clairement imputable a un retournement de tendance non anticipe du MACD, et non a une defaillance fondamentale des algorithmes.

\end{document}
